% Part: incompleteness
% Chapter: theories-computability
% Section: extensions-of-q-not-decidable

\documentclass[../../../include/open-logic-section]{subfiles}

\begin{document}

\olfileid{inc}{tcp}{cqn}

\olsection{Ekstensi $\Th{Q}$ kang konsisten ora bisa diputusake}

\begin{explain}
Elinga yen teori iku \emph{konsisten} yen ora mbuktekake $!A$ lan
$\lnot !A$ bebarengan kanggo formula~$!A$ sembarang. Amarga apa wae
bisa diturunake saka kontradiksi, teori kang ora konsisten iku trivial:
saben ukara bisa dibuktekake. Cetha yen teori iku $\omega$-konsisten,
mula teori iku konsisten. Nanging konsisten iku syarat kang luwih
ringkih (yaiku, ana teori kang konsisten nanging ora $\omega$-konsisten).
Kita bisa ngendhori panganggep ing \olref[oqn]{thm:oconsis-q} dadi
kekonsistenan lumrah kanggo ngasilake teorema kang luwih kuwat.
\end{explain}

\begin{lem}
Ora ana ``relasi komputabel universal''. Tegese, ora ana relasi
komputabel biner $R(x,y)$ kang nduweni sipat mangkene:
kapan wae $S(y)$ iku relasi komputabel uner, ana sawijining $k$ saengga
kanggo saben $y$, $S(y)$ bener yen lan mung yen $R(k,y)$ bener.
\end{lem}

\begin{proof}
Upamana $R(x,y)$ iku relasi komputabel universal. Upamana $S(y)$
iku relasi $\lnot R(y,y)$. Amarga $S(y)$ komputabel, kanggo sawijining
$k$, $S(y)$ ekuivalen karo $R(k,y)$. Nanging banjur $S(k)$
ekuivalen karo $R(k,k)$ lan $\lnot R(k,k)$ bebarengan, kang dadi
kontradiksi.
\end{proof}


\begin{thm}
Upamana $\Th{T}$ teori konsisten sembarang kang ngemot $\Th{Q}$. Mula
$\Th{T}$ ora bisa diputusake.
\end{thm}


\begin{proof}
Upamana $\Th{T}$ iku ekstensi $\Th{Q}$ kang konsisten lan bisa
diputusake. Kita bakal nemokake kontradiksi kanthi nggunakake $\Th{T}$
kanggo nemtokake relasi komputabel universal.

Tetepna $R(x,y)$ bener yen lan mung yen
\begin{quote}
$x$ ngodeake formula $!D(u)$, lan $\Th{T}$ mbuktekake $!D(\num y)$.
\end{quote}
Amarga kita nganggep yen $\Th{T}$ bisa diputusake, $R$ iku komputabel.
Saiki ayo dituduhake yen $R$ universal. Yen $S(y)$ relasi komputabel
sembarang, mula relasi iki bisa direpresentasikake ing $\Th{Q}$
(lan mulane $\Th{T}$) dening formula $!D_S(u)$. Mula kanggo saben $n$,
kita nduweni
\begin{eqnarray*}
S(\num n) & \lif & T \vdash !D_S(\num n) \\
& \lif & R(\Gn{!D_S(u)}, n)
\end{eqnarray*}
and
\begin{eqnarray*}
\lnot S(\num n) & \lif & T \vdash \lnot !D_S(\num n) \\
& \lif & T \not\vdash !D_S(\num n) \quad \text{(amarga $\Th{T}$
  konsisten)} \\
& \lif & \lnot R(\Gn{!D_S(u)}, n).
\end{eqnarray*}
Tegese, kanggo saben $y$, $S(y)$ bener yen lan mung yen
$R(\Gn{!D_S(u)}, y)$ bener. Mula $R$ universal, lan kita nemokake
kontradiksi kang digoleki.
\end{proof}

Upamana ``aritmetika bener'' iku teori $\Setabs{!A}{ \Sat{\Nat}{!A}}$,
yaiku himpunan ukara ing basa aritmetika kang bener ing interpretasi
standar.

\begin{cor}
Aritmetika bener ora bisa diputusake.
\end{cor}

%In Section~\olref{undefinability:truth:section} we will state a stronger
%result, due to Tarski.

\end{document}

